\section{Transformacje naturalne}\label{sec:naturalne}

\providecommand{\atlasref}[2]{#2}

Dwa funktory mogą przyporządkowywać obiektom różne struktury, a mimo to być ze sobą systematycznie powiązane. Pojedyncza rodzina odwzorowań nie wystarczy: trzeba jeszcze wymagać zgodności z każdym morfizmem kategorii wyjściowej.

\begin{definicja}[Transformacja naturalna]\label{def:naturalna}
Dla funktorów $F,G:\mathcal C\to\mathcal D$ \emph{transformacja naturalna} $\eta:F\Rightarrow G$ jest rodziną morfizmów $\eta_A:F(A)\to G(A)$, po jednym dla każdego $A\in\Ob(\mathcal C)$, spełniającą dla każdego $f:A\to B$ warunek $G(f)\circ\eta_A=\eta_B\circ F(f)$. Innymi słowy, poniższy diagram jest przemienny:
\[
\begin{CD}
F(A) @>{\eta_A}>> G(A)\\
@V{F(f)}VV @VV{G(f)}V\\
F(B) @>>{\eta_B}> G(B).
\end{CD}
\]
Jeśli wszystkie $\eta_A$ są izomorfizmami, to $\eta$ nazywamy \emph{izomorfizmem naturalnym} i piszemy $F\cong G$.
\end{definicja}

\begin{stwierdzenie}[Funktory i transformacje naturalne między posetami]\label{tw:fun-posety}
Niech $(P,\le_P)$ i $(Q,\le_Q)$ będą zbiorami częściowo uporządkowanymi, traktowanymi jako kategorie. Funktory $F:P\to Q$ są dokładnie funkcjami monotonicznymi. Jeśli $F,G:P\to Q$ są monotoniczne, to istnieje transformacja naturalna $F\Rightarrow G$ wtedy i tylko wtedy, gdy $F(x)\le_Q G(x)$ dla każdego $x\in P$; jest ona wówczas jednoznaczna.
\end{stwierdzenie}
\begin{proof}
W $P$ istnieje morfizm $x\to y$ dokładnie wtedy, gdy $x\le_P y$. Funktor musi więc wysyłać każdą taką parę w parę $F(x)\le_Q F(y)$, co jest monotonicznością. Odwrotnie, funkcja monotoniczna definiuje morfizm $F(x)\to F(y)$ dla każdego $x\to y$; identyczności i złożenia są zachowane, gdyż odpowiednie zbiory morfizmów w $Q$ mają co najwyżej jeden element. Komponent transformacji naturalnej $\eta_x:F(x)\to G(x)$ istnieje dokładnie przy wskazanej nierówności. Dla $x\le_P y$ obie strony równania naturalności są wówczas morfizmami $F(x)\to G(y)$, a więc są równe.
\end{proof}

\begin{przyklad}[Zbiory uporządkowane przez inkluzję]
W kategorii pochodzącej ze zbioru częściowo uporządkowanego $(\mathcal P(\{1,2,3\}),\subseteq)$ istnieje strzałka $\{1\}\to\{1,2\}$, ale nie strzałka odwrotna. Funkcja $F(U)=U\cap\{1,2\}$ jest monotoniczna, zatem daje funktor. Ponadto $F(U)\subseteq U$ dla każdego $U$, co określa transformację naturalną $F\Rightarrow\id$. Nie jest ona izomorfizmem naturalnym, bo np. $F(\{3\})=\varnothing\neq\{3\}$.
\end{przyklad}

\begin{przyklad}[Naturalne odwzorowanie do przestrzeni podwójnie dualnej]\label{przyklad:podwojny-dual}
W $\Vect_k$ dla każdej przestrzeni $V$ definiujemy $\iota_V:V\to V^{**}$ przez $\iota_V(v)(\lambda)=\lambda(v)$ dla $v\in V$, $\lambda\in V^*$. Jest to odwzorowanie liniowe. Jeśli $f:V\to W$, to $f^{**}:V^{**}\to W^{**}$ spełnia
$f^{**}(\iota_V(v))(\mu)=\iota_V(v)(f^*(\mu))=\mu(f(v))=\iota_W(f(v))(\mu)$
dla wszystkich $v\in V$ i $\mu\in W^*$. Stąd $f^{**}\circ\iota_V=\iota_W\circ f$: rodzina $(\iota_V)$ jest transformacją naturalną $\id_{\Vect_k}\Rightarrow(-)^{**}$.

Odwzorowanie $\iota_V$ jest różnowartościowe: jeżeli $v\ne0$, rozszerzamy $\{v\}$ do bazy Hamela i definiujemy funkcjonał $\lambda$ przyjmujący wartość $1$ na $v$ i $0$ na pozostałych wektorach bazy. Wtedy $\iota_V(v)(\lambda)=1$. Używamy tutaj twierdzenia o istnieniu bazy (w ZFC). Jeżeli $\dim V=n<\infty$, z bazy $V$ otrzymujemy bazę dualną, więc $\dim V^{**}=\dim V^*=n$. Różnowartościowe odwzorowanie między przestrzeniami tej samej skończonej wymiarowości jest bijekcją. Na $\Vect_k^{\mathrm{fd}}$ otrzymujemy zatem izomorfizm naturalny.

Dla przykładu niech $V=k^{(\mathbb N)}$ będzie przestrzenią ciągów o skończonym nośniku z bazą $(e_j)_{j\in\mathbb N}$. Jej przestrzeń dualna utożsamia się z $k^{\mathbb N}$ przez wartości funkcjonału na $e_j$. W ilorazie $k^{\mathbb N}/k^{(\mathbb N)}$ klasa stałego ciągu $(1,1,\ldots)$ jest niezerowa, więc istnieje funkcjonał przyjmujący na niej wartość $1$ (ponownie przez rozszerzenie do bazy). Złożenie tego funkcjonału z rzutem ilorazowym daje $\Lambda\in V^{**}$, który zeruje wszystkie ciągi o skończonym nośniku, ale nie jest zerowy. Gdyby $\Lambda=\iota_V(v)$, podstawienie ciągu mającego jedynkę na miejscu $j$ i zera na pozostałych miejscach dawałoby $v_j=0$ dla każdego $j$. Wtedy $v=0$, sprzecznie z $\Lambda\ne0$. Zatem $\iota_V$ nie jest surjektywna w tym przykładzie.
\end{przyklad}

\begin{przyklad}[Naturalne przyporządkowanie singletonu]
Dla kowariantnego funktora potęgowego $\mathcal P_+$ z sekcji~\ref{sec:funktory} definiujemy $\eta_X:X\to\mathcal P(X)$ przez $\eta_X(x)=\{x\}$. Dla funkcji $f:X\to Y$ mamy $\mathcal P_+(f)(\eta_X(x))=\{f(x)\}=\eta_Y(f(x))$. Zatem $\eta:\id_{\Set}\Rightarrow\mathcal P_+$ jest naturalna.
\end{przyklad}

\begin{przyklad}[Naturalność i działania grup]
Niech $G$ będzie grupą, a $F,H:G\to\Set$ funktorami reprezentującymi lewe działania $G$ na zbiorach $X$ i $Y$ (przykład~\ref{przyklad:dzialanie-funktor}). Transformacja naturalna $\eta:F\Rightarrow H$ składa się z jednej funkcji $\eta_*:X\to Y$. Warunek naturalności dla morfizmu $g\in G$ ma postać $\eta_*(g\cdot x)=g\cdot\eta_*(x)$. Jest to dokładnie definicja odwzorowania $G$-ekwiwariantnego.
\end{przyklad}

\begin{twierdzenie}[Naturalne endomorfizmy identyczności w przestrzeniach liniowych]\label{tw:skalarnosc}
Niech $k$ będzie ciałem. Każda transformacja naturalna $\eta:\id_{\Vect_k}\Rightarrow\id_{\Vect_k}$ ma postać $\eta_V=\lambda\id_V$ dla pewnego $\lambda\in k$, niezależnego od $V$. Odwrotnie, każdy $\lambda\in k$ definiuje taką transformację naturalną.
\end{twierdzenie}
\begin{proof}
Ponieważ $\eta_k:k\to k$ jest liniowa, istnieje $\lambda=\eta_k(1)\in k$ takie, że $\eta_k(t)=\lambda t$. Dla dowolnego $V$ i $v\in V$ rozważmy liniową funkcję $f_v:k\to V$, $t\mapsto tv$. Naturalność daje $\eta_V(f_v(1))=f_v(\eta_k(1))$, więc $\eta_V(v)=\lambda v$. Odwrotnie, dla dowolnego liniowego $f:V\to W$ mamy $f(\lambda v)=\lambda f(v)$, co jest warunkiem naturalności.
\end{proof}

\begin{uwaga}[Naturalny nie znaczy dowolnie wybrany]
Każda skończenie wymiarowa przestrzeń wektorowa ma bazę, ale \emph{wybór bazy} zazwyczaj nie jest naturalny: różne automorfizmy tej samej przestrzeni nie muszą zachowywać wybranej bazy. Natomiast odwzorowanie $V\to V^{**}$ z przykładu~\ref{przyklad:podwojny-dual} nie zależy od wyboru bazy. Naturalność formalizuje to rozróżnienie.
\end{uwaga}

\begin{stwierdzenie}[Składanie pionowe transformacji]\label{cat:prop:vertical}
Jeśli $F,G,H:\mathcal C\to\mathcal D$ i $\eta:F\Rightarrow G$, $\theta:G\Rightarrow H$ są transformacjami naturalnymi, to rodzina $(\theta\circ\eta)_A=\theta_A\circ\eta_A$ jest transformacją naturalną $F\Rightarrow H$. Dla małych $\mathcal C$ i lokalnie małych $\mathcal D$ funktory $\mathcal C\to\mathcal D$ z transformacjami naturalnymi tworzą kategorię $[\mathcal C,\mathcal D]$.
\end{stwierdzenie}
\begin{proof}
Dla $f:A\to B$ mamy $H(f)\theta_A\eta_A=\theta_B G(f)\eta_A=\theta_B\eta_B F(f)$, gdzie pominięto znaki składania. Łączność wynika z łączności w $\mathcal D$ w każdej składowej, a identycznością funktora $F$ jest rodzina $(\id_{F(A)})_A$. Jest ona naturalna, gdyż $F(f)\circ\id_{F(A)}=\id_{F(B)}\circ F(f)$; złożenie z nią nie zmienia żadnej składowej. Jeżeli $\mathcal C$ jest mała, zbiór transformacji $F\Rightarrow G$ jest podzbiorem produktu zbiorów $\prod_{A\in\Ob(\mathcal C)}\Hom_{\mathcal D}(F(A),G(A))$. Produkt jest zbiorem, ponieważ indeksujemy go zbiorem obiektów i $\mathcal D$ jest lokalnie mała. Kategoria funktorów jest więc lokalnie mała. Nie twierdzimy, że jest mała, jeśli $\mathcal D$ jest duża.
\end{proof}


\begin{twierdzenie}[Jedyna naturalna operacja jednoargumentowa na zbiorach]\label{tw:only-id-set}
Każda transformacja naturalna $\eta:\id_{\Set}\Rightarrow\id_{\Set}$ jest identycznościowa.
\end{twierdzenie}
\begin{proof}
Dla jednoelementowego zbioru $1=\{*\}$ istnieje tylko jedna funkcja $1\to1$, zatem $\eta_1=\id_1$. Niech $X$ będzie dowolnym zbiorem, a $x\in X$. Dla funkcji $f_x:1\to X$, $f_x(*)=x$, naturalność daje $\eta_X(x)=\eta_X(f_x(*))=f_x(\eta_1(*))=x$. Ponieważ równość zachodzi dla każdego $x\in X$, otrzymujemy $\eta_X=\id_X$ (również gdy $X=\varnothing$, gdzie funkcja $X\to X$ jest jednoznaczna).
\end{proof}

\begin{przyklad}[Dlaczego wybór bazy nie jest naturalny]
Dla $V=\mathbb R^2$ wybierzmy bazę $(v_1,v_2)$. Istnieje automorfizm liniowy $T:V\to V$ spełniający $T(v_1)=v_1+v_2$ i $T(v_2)=v_2$. Wybór uporządkowanej bazy zgodny z \emph{każdym} izomorfizmem wymagałby, aby automorfizm $T$ zachowywał oba wybrane wektory, co jest fałszem. Przeciwieństwem jest naturalne odwzorowanie $V\to V^{**}$: dla niego zgodność z dowolnym $T$ została wykazana w przykładzie~\ref{przyklad:podwojny-dual}.
\end{przyklad}

\begin{stwierdzenie}[Kiedy transformacja naturalna jest odwracalna]\label{tw:nat-iso-components}
Niech $F,G:\mathcal C\to\mathcal D$ będą funktorami i $\eta:F\Rightarrow G$ transformacją naturalną. Istnieje transformacja naturalna $\theta:G\Rightarrow F$ z $\theta\circ\eta=\id_F$ i $\eta\circ\theta=\id_G$ wtedy i tylko wtedy, gdy każda składowa $\eta_A:F(A)\to G(A)$ jest izomorfizmem w $\mathcal D$.
\end{stwierdzenie}
\begin{proof}
Jeśli $\theta$ jest odwrotnością naturalną, to $\theta_A\eta_A=\id_{F(A)}$ i $\eta_A\theta_A=\id_{G(A)}$ dla każdego $A$, więc $\eta_A$ jest izomorfizmem. Odwrotnie, załóżmy odwracalność wszystkich $\eta_A$ i zdefiniujmy $\theta_A=\eta_A^{-1}$. Dla $f:A\to B$ naturalność $\eta$ to $G(f)\eta_A=\eta_BF(f)$. Składając z lewej strony z $\eta_B^{-1}$, a z prawej z $\eta_A^{-1}$, otrzymujemy $F(f)\eta_A^{-1}=\eta_B^{-1}G(f)$, czyli naturalność $\theta$. Równości $\theta\eta=\id_F$ i $\eta\theta=\id_G$ zachodzą składowa po składowej.
\end{proof}

\subsection{Składanie poziome i prawo wymiany}

\begin{definicja}[Składanie poziome]\label{cat:def:horizontal}
Niech $F,G:\mathcal C\to\mathcal D$, $H,K:\mathcal D\to\mathcal E$ oraz $\alpha:F\Rightarrow G$, $\beta:H\Rightarrow K$. Składanie pionowe zapisujemy dalej jako $\circ$, a składanie poziome jako $\ast$. Definiujemy
\[
(\beta\ast\alpha)_A=\beta_{G(A)}\circ H(\alpha_A)
=K(\alpha_A)\circ\beta_{F(A)}:HF(A)\to KG(A).
\]
Zapis $HF$ oznacza złożenie funktorów $H\circ F$. Równość dwóch wzorów wynika z naturalności $\beta$ względem morfizmu $\alpha_A$ w $\mathcal D$. Rodzinę $\beta\ast\alpha:HF\Rightarrow KG$ nazywamy złożeniem poziomym.
\end{definicja}

\begin{stwierdzenie}\label{cat:prop:horizontal}
Złożenie poziome jest transformacją naturalną. W szczególności można składać transformację naturalną z funktorem z lewej lub z prawej strony.
\end{stwierdzenie}
\begin{proof}
Dla $f:A\to B$ korzystamy kolejno z naturalności $\beta$ względem $G(f)$ i naturalności $\alpha$:
\begin{align*}
KG(f)\circ(\beta\ast\alpha)_A
&=K(G(f))\circ\beta_{G(A)}\circ H(\alpha_A)\\
&=\beta_{G(B)}\circ H(G(f))\circ H(\alpha_A)\\
&=\beta_{G(B)}\circ H(G(f)\circ\alpha_A)\\
&=\beta_{G(B)}\circ H(\alpha_B\circ F(f))\\
&=(\beta\ast\alpha)_B\circ HF(f).
\end{align*}
Jeżeli $\beta=\id_H$, otrzymujemy transformację $H\alpha:HF\Rightarrow HG$ o składowych $H(\alpha_A)$. Jeżeli $\alpha=\id_F$, otrzymujemy $\beta F:HF\Rightarrow KF$ o składowych $\beta_{F(A)}$.
\end{proof}

\begin{twierdzenie}[Prawo wymiany]\label{cat:thm:interchange}
Niech $F_0,F_1,F_2:\mathcal C\to\mathcal D$, $G_0,G_1,G_2:\mathcal D\to\mathcal E$ oraz
\[
\alpha:F_0\Rightarrow F_1,\quad \alpha':F_1\Rightarrow F_2,
\qquad
\beta:G_0\Rightarrow G_1,\quad \beta':G_1\Rightarrow G_2.
\]
Wtedy jako transformacje $G_0F_0\Rightarrow G_2F_2$ mamy
\[
(\beta'\circ\beta)\ast(\alpha'\circ\alpha)
=(\beta'\ast\alpha')\circ(\beta\ast\alpha).
\]
\end{twierdzenie}
\begin{proof}
W obiekcie $A$ prawa strona jest równa
\[
\beta'_{F_2(A)}\circ G_1(\alpha'_A)\circ
\beta_{F_1(A)}\circ G_0(\alpha_A).
\]
Naturalność $\beta$ względem $\alpha'_A:F_1(A)\to F_2(A)$ daje
$G_1(\alpha'_A)\circ\beta_{F_1(A)}
=\beta_{F_2(A)}\circ G_0(\alpha'_A)$.
Po podstawieniu i wykorzystaniu funktorialności $G_0$ otrzymujemy
\[
(\beta'_{F_2(A)}\circ\beta_{F_2(A)})\circ
G_0(\alpha'_A\circ\alpha_A),
\]
czyli składową lewej strony. Równość wszystkich składowych dowodzi tezy.
\end{proof}

\begin{stwierdzenie}[Łączność i jednostki składania poziomego]\label{cat:prop:horizontal-assoc}
Dla dodatkowych funktorów $L,M:\mathcal E\to\mathcal K$ i transformacji $\gamma:L\Rightarrow M$ zachodzi
$\gamma\ast(\beta\ast\alpha)=(\gamma\ast\beta)\ast\alpha$.
Składanie poziome z transformacją identycznościową funktora identycznościowego odpowiedniej kategorii nie zmienia transformacji.
\end{stwierdzenie}
\begin{proof}
W obiekcie $A$ lewa strona ma składową
\[
\gamma_{KG(A)}\circ L(\beta_{G(A)}\circ H(\alpha_A))
=\gamma_{KG(A)}\circ L(\beta_{G(A)})\circ LH(\alpha_A).
\]
Prawa strona ma tę samą składową z definicji składania poziomego. Dla jednostki z lewej strony wzór daje $\id_{G(A)}\circ\alpha_A=\alpha_A$. Dla jednostki z prawej strony zastosowany do $\alpha\ast\id_{\id_{\mathcal C}}$ daje $\alpha_A\circ F(\id_A)=\alpha_A$.
\end{proof}

\subsection{Naturalność konstrukcji algebraicznych i geometrycznych}

\begin{przyklad}[Naturalny izomorfizm tensorowy]\label{cat:ex:tensor-unit-natural}
Dla pierścienia przemiennego $R$ z jedynką odwzorowanie $\ell_N:R\otimes_R N\to N$, $r\otimes n\mapsto rn$, jest izomorfizmem o odwrotności $n\mapsto1\otimes n$. Istnienie i wzajemną odwrotność wykazano w rozdziale 18 skryptu \emph{Algebra abstrakcyjna}. Dla homomorfizmu $f:N\to P$ mamy
\[
f(\ell_N(r\otimes n))=f(rn)=rf(n)=\ell_P((\id_R\otimes f)(r\otimes n)).
\]
Równość na tensorach prostych wystarcza, bo generują one moduł. Otrzymujemy izomorfizm naturalny $R\otimes_R(-)\cong\id$ w kategorii $R$-modułów.
\end{przyklad}

\begin{przyklad}[Naturalność odwzorowania łączącego]\label{cat:ex:connecting-natural}
Rozważmy kategorię krótkich ciągów dokładnych kompleksów łańcuchowych $0\to A_\bullet\xrightarrow{i}B_\bullet\xrightarrow{p}C_\bullet\to0$ nad ustalonym pierścieniem. Morfizmy są trójkami odwzorowań łańcuchowych $(f_A,f_B,f_C)$ zgodnymi z $i,p$, czyli $f_Bi=i'f_A$ oraz $p'f_B=f_Cp$. Składamy je współrzędnie. Funktory przypisujące ciągowi $H_n(C_\bullet)$ oraz $H_{n-1}(A_\bullet)$ działają na morfizmach przez $H_n(f_C)$ oraz $H_{n-1}(f_A)$.

Odwzorowania łączące $\delta_n:H_n(C_\bullet)\to H_{n-1}(A_\bullet)$, skonstruowane w \atlasref{TOP-03::uwaga:les-naturalnosc}{TOP-03: naturalność odwzorowania łączącego}, tworzą transformację naturalną tych funktorów. Rzeczywiście, dla cyklu $c\in C_n$ wybieramy $b\in B_n$ z $p(b)=c$. Ponieważ $p(d_Bb)=d_Cc=0$, istnieje jedyne $a\in A_{n-1}$ z $i(a)=d_Bb$. Jest cyklem, gdyż $i(d_Aa)=d_B^2b=0$ i $i$ jest różnowartościowe. Definicja odwzorowania łączącego daje $\delta_n[c]=[a]$.

Dla drugiego ciągu $f_B(b)$ jest podniesieniem $f_C(c)$, gdyż $p'f_B(b)=f_Cp(b)$. Ponadto
\[
d_{B'}f_B(b)=f_B(d_Bb)=f_Bi(a)=i'f_A(a).
\]
A zatem $\delta'_n(H_n(f_C)[c])=[f_A(a)]=H_{n-1}(f_A)(\delta_n[c])$. Jest to warunek naturalności. Poprawność samej konstrukcji $\delta_n$ na klasach ilorazowych pochodzi z przywołanego rozdziału.
\end{przyklad}

\begin{przyklad}[Izomorfizm naturalny de Rhama]\label{cat:ex:derham-natural}
Twierdzenie de Rhama (udowodnione w rozdziale 10 skryptu \emph{Geometria i topologia różniczkowa}) daje izomorfizmy
\[
I_M:H^q_{\mathrm{dR}}(M)\longrightarrow H^q_{\mathrm{sing}}(M;\mathbb R).
\]
Sprawdźmy sens ich naturalności. Na kołańcuchach określonych na gładkich sympleksach odwzorowanie całkowania ma postać $I_M(\omega)(\sigma)=\int_{\Delta^q}\sigma^*\omega$. Jeżeli $f:M\to N$ jest gładkie, to dla formy $\omega$ na $N$ i gładkiego sympleksu $\sigma:\Delta^q\to M$ mamy
\[
I_M(f^*\omega)(\sigma)=\int_{\Delta^q}(f\circ\sigma)^*\omega
=I_N(\omega)(f\circ\sigma).
\]
Prawa strona jest wartością cofnięcia kołańcucha $I_N(\omega)$ na $\sigma$. Twierdzenie Stokesa zapewnia zgodność całkowania z różniczkami kompleksów. Ograniczenie kołańcuchów singularnych do gładkich sympleksów także komutuje z cofnięciami; w przywołanym dowodzie indukuje ono izomorfizm na kohomologii. Jego odwrotność jest naturalna na mocy stwierdzenia~\ref{tw:nat-iso-components}. W konsekwencji $I_M\circ H^q_{\mathrm{dR}}(f)=H^q_{\mathrm{sing}}(f)\circ I_N$, czyli otrzymujemy izomorfizm dwóch funktorów kontrawariantnych.
\end{przyklad}

\begin{przyklad}[Zapowiedź porównania GAGA]\label{cat:ex:gaga-natural}
Niech $X$ będzie schematem właściwym skończonego typu nad $\mathbb C$. Twierdzenie Serre'a GAGA, przywołane jako wynik zewnętrzny w rozdziale 9 skryptu \emph{Geometria algebraiczna}, daje dla ustalonego $q\ge0$ naturalne izomorfizmy
\[
H^q(X,\mathcal F)\longrightarrow H^q(X^{\mathrm{an}},\mathcal F^{\mathrm{an}})
\]
dla koherentnych snopów algebraicznych $\mathcal F$ na $X$. Naturalność dotyczy morfizmów takich snopów: porównanie kohomologii komutuje z odwzorowaniami indukowanymi przez morfizm $\mathcal F\to\mathcal G$ i jego analityfikację. Jest to zapowiedź zastosowania transformacji naturalnych do snopów; nie dowodzimy tu twierdzenia GAGA ani nie usuwamy założenia właściwości.
\end{przyklad}

